% from ega3-7-fr.tex
\oldpage[III]{186}
est un foncteur homologique de $\mathbf{Ab}_A$ dans $\mathbf{Ab}$,
évidemment $A$-linéaire lorsque $A$ est commutatif (7.1.2), et commutant aux
limites inductives.

Si $A$ est commutatif, alors, pour toute $A$-algèbre $B$, le foncteur
homologique $T_\bullet^{(B)}$ (7.3.8) est donné par définition par
\[
  T_\bullet^{(B)}(N)
  =
  H_\bullet(P_\bullet\otimes_A N_{[\rho]})
  \tag{7.4.1.2}\label{III.7.4.1.2.fr}
\]
où $\rho:A\to B$ est l'homomorphisme définissant la structure d'algèbre de
$B$; comme on peut aussi écrire
\[
  P_\bullet\otimes_A N_{[\rho]}
  =P_\bullet\otimes_A(B\otimes_B N)_{[\rho]}
  =(P_\bullet\otimes_A B)\otimes_B N,
\]
on voit que l'on a
\[
  T_\bullet^{(B)}(N)=H_\bullet(P'_\bullet\otimes_B N)
  \tag{7.4.1.3}\label{III.7.4.1.3.fr}
\]
pour tout $B$-module $N$, $P'_\bullet$ étant le complexe
$P_\bullet\otimes_A B$ de $B$-modules plats (0\textsubscript{I}, 6.2.1).
\end{env}

\begin{proposition}[7.4.2]
Sous les conditions générales de (7.4.1), et pour un entier $p\in\mathbf Z$
donné, les propriétés suivantes sont équivalentes :
\begin{enumerate}
  \item[a)] $T_p$ est exact à gauche (ou, ce qui revient au même, $T_{p+1}$
    est exact à droite).
  \item[b)]
    \[
      Z'_p(P_\bullet)
      =
      \operatorname{Coker}(P_{p+1}\longrightarrow P_p)
    \]
    est un $A$-module à droite plat.
  \item[c)] Il existe un complexe $P'_\bullet$ de $A$-modules à droite plats
    tel que la différentielle
    \[
      d'_{p+1}:P'_{p+1}\longrightarrow P'_p
    \]
    soit nulle, et un isomorphisme de foncteurs homologiques de
    $H_\bullet(P_\bullet\otimes_A M)$ sur
    $H_\bullet(P'_\bullet\otimes_A M)$.
\end{enumerate}
\end{proposition}

\begin{proof}
Par définition, on a une suite exacte fonctorielle en $M$
\[
  0\longrightarrow T_p(M)
  \longrightarrow Z'_p(P_\bullet\otimes M)
  \longrightarrow P_{p-1}\otimes M
\]
où
\[
  Z'_p(P_\bullet\otimes M)
  =
  \operatorname{Coker}(P_{p+1}\otimes M\longrightarrow P_p\otimes M)
  =
  Z'_p(P_\bullet)\otimes M
\]
en vertu de l'exactitude à droite du produit tensoriel. Pour tout homomorphisme
$f:M\to N$, on a donc un diagramme commutatif
\[
\xymatrix@C=30pt@R=38pt{
  0\ar[r]
  &T_p(M)\ar[r]\ar[d]_u
  &Z'_p(P_\bullet)\otimes M\ar[r]\ar[d]_v
  &P_{p-1}\otimes M\ar[d]_w
  \\
  0\ar[r]
  &T_p(N)\ar[r]
  &Z'_p(P_\bullet)\otimes N\ar[r]
  &P_{p-1}\otimes N
}
\tag{7.4.2.1}\label{III.7.4.2.1.fr}
\]
dont les lignes sont exactes. Si $f$ est un monomorphisme, il en est de même
de $w$ puisque $P_{p-1}$ est plat; si $T_p$ est exact à gauche, $u$ est lui
aussi un monomorphisme; on en conclut que $v$ est un monomorphisme, ce qui
entraîne que $Z'_p(P_\bullet)$ est plat. Inversement, s'il en est ainsi, $v$
est un monomorphisme pour tout monomorphisme $f:M\to N$, donc le diagramme
(7.4.2.1) montre que $u$ est un monomorphisme, et par suite $T_p$ (qui est déjà
semi-exact) est exact à gauche. Ainsi $a)$ et $b)$ sont équivalentes. Il est
immédiat que $c)$ entraîne $a)$, car si
$d'_{p+1}:P'_{p+1}\to P'_p$ est nulle, et
$0\to M'\to M\to M''\to0$ est une suite exacte de $A$-modules, l'opérateur
bord dans la suite exacte
\[
  H_{p+1}(P'_\bullet\otimes M'')
  \xrightarrow{\partial}H_p(P'_\bullet\otimes M')
  \longrightarrow H_p(P'_\bullet\otimes M)
\]
